\section{The statement must be right before the proof can help}\label{ch16:sec:statement}
When Lean accepts a proof, it guarantees one thing: the conclusion, exactly as it is written in Lean, follows from the hypotheses, exactly as they are written in Lean. It does not guarantee that the Lean statement says what anyone meant. A Lean statement is a translation of a mathematical sentence, perhaps from a research paper, a textbook exercise, or a request to an AI assistant, and a translation can be wrong even when the proof of it is perfectly correct. Lean cannot notice such a mistake, because it never sees the original sentence. Comparing the formal statement with the intended one is a separate task, and it is the reader's job, not the proof assistant's~\cite{leanvalidate}.

Here is how easily the two can drift apart. Section~\ref{ch07:sec:seymour} defined the second out-neighborhood $N_2^+(v)$ of a vertex $v$ as the set of vertices at directed distance exactly two from $v$. A tempting shortcut is to take instead all the endpoints $w$ of directed walks $v\to u\to w$ of length two. Both descriptions are precise, and Lean would accept correct proofs about either one. But they describe different sets. Take the oriented graph from that section, with edges $a\to b$, $b\to c$, $a\to c$, and $c\to d$. The two-step walks from $a$ are $a\to b\to c$ and $a\to c\to d$, so the shortcut gives the set $\{c,d\}$. The only vertex at directed distance exactly two from $a$ is $d$, because the edge $a\to c$ puts $c$ at distance one. Since the shortcut can only add vertices to the second out-neighborhood, a ``proof of Seymour's conjecture'' that used it would prove a weaker statement than the conjecture. A Lean definition records every such choice, but someone still has to read the definition and notice which choice was made.

So before reading any tactics, compare the formal statement with the intended one, piece by piece. Make a table with two columns. In the first, list the objects of the mathematical statement, its hypotheses, and its conclusion; in the second, write what represents each of them in Lean. This is the statement contract of Habit~2 (Chapter~\ref{ch:logic}), written out for a formal statement, and Section~\ref{ch16:sec:surj} gives a complete example. While you fill in the table, check the places where translations most often go wrong. Seymour's conjecture, \emph{every nonempty finite oriented graph has a vertex $v$ with $|N_2^+(v)|\ge|N_1^+(v)|$}, shows what can happen at each of them.
\begin{itemize}
\item \emph{Domains.} Does each variable range over the intended set? In Section~\ref{ch15:sec:exists}, the expression $7-10$ meant $0$ for natural numbers and $-3$ for integers.
\item \emph{Strict and non-strict inequalities.} With $>$ in place of $\ge$, the conjecture would be false. A graph with one vertex and no edges is a nonempty finite oriented graph. Its only vertex $v$ has $|N_2^+(v)|=0=|N_1^+(v)|$, so no vertex satisfies $|N_2^+(v)|>|N_1^+(v)|$.
\item \emph{Nonemptiness.} A graph with no vertices has no vertex with the required property, simply because it has no vertex at all. Without the word ``nonempty,'' the conjecture would be false.
\item \emph{Excluded cases.} The conditions $w\ne v$ and $w\notin N_1^+(v)$ in the definition of $N_2^+(v)$ are exactly what the shortcut above leaves out.
\item \emph{The order of quantifiers.} The conjecture claims that \emph{some} vertex satisfies the inequality, not that every vertex does. In the same way, ``for every $y\in B$ there is an $x\in A$ with $f(x)=y$'' says that $f$ is surjective, whereas ``there is an $x\in A$ with $f(x)=y$ for every $y\in B$'' asks a single input to be sent to every element of $B$ at once. That is impossible as soon as $B$ has two different elements (Section~\ref{ch02:sec:quantifiers}).
\end{itemize}

\subsection*{What to read first in an unfamiliar theorem}
When you meet an unfamiliar Lean theorem, read it in this order. First find the whole declaration, from the word \code{theorem} to the \code{:=} that begins the proof. Its parameters name the objects and the hypotheses, and the proposition after the colon that follows the parameters is the conclusion (Section~\ref{ch15:sec:symbols}). Next, look up each definition that appears in the statement and read its \emph{body}, the part after its own \code{:=}. A definition can contain conditions that its short name does not reveal. Only when you understand the statement should you turn to the proof and see how it builds its evidence.

Names alone are not enough. The name \code{Second} in Section~\ref{ch16:sec:second} does not say whether first out-neighbors are excluded; only the body of the definition does. The name \code{Surj} in the next section does not say which set is the domain and which is the codomain, or in which order the quantifiers come. Section~\ref{ch15:sec:symbols} made the same point with a theorem that could be called \code{all\_problems\_solved} while stating only \code{0 = 0}.
